\section*{अध्याय \ref{ch:b3:quantitative-geomorphology} --- परिमाणात्मक भू-आकृति विज्ञान}

\begin{solution}{exo:b3:quantitative-geomorphology:1}
$a/\sqrt{\pi\kappa t} = 0.30 - 0.04 = 0.26$, इसलिए $\kappa t = (2/0.26)^2/\pi =
\qty{18.8}{m^2}$, और $t = 18.8/(5\times10^{-3}) = \qty{3.8}{kyr}$, दो गुने के भीतर
दोनों ओर, $\kappa$ के परास को देखते हुए।
\end{solution}

\begin{solution}{exo:b3:quantitative-geomorphology:2}
$E = \kappa\,|\nabla^2 z| = 4\times10^{-3}\times0.012 =
\qty{4.8e-5}{m/yr}$, \qty{0.048}{mm/yr}। उच्चावच $E L^2/(2\kappa) =
0.012\times80^2/2 = \qty{38.4}{m}$।
\end{solution}

\begin{solution}{exo:b3:quantitative-geomorphology:3}
$E = 2\times10^{-6}\times(25\times10^6)^{0.5}\times0.01 = 2\times10^{-6}
\times5000\times0.01 = \qty{1.0e-4}{m/yr}$, \qty{0.1}{mm/yr}। चूँकि $E$ का
मात्रक \unit{m/yr} है, $A^m$ का $\unit{m}^{2m}$ है, $S$ बेमात्रक है, और $K$ का
$\unit{m}^{1-2m}\,\unit{yr^{-1}}$ मात्रक है: $m = 0.5$ के लिए \unit{yr^{-1}}।
\end{solution}

\begin{solution}{exo:b3:quantitative-geomorphology:4}
(क) स्थिर: एक कठोर संस्तर जहाँ उभरता है वहाँ नदी को तीव्र कर देता है, और सोपान
संस्तर के साथ ही बना रहता है। (ख) भ्रंश पर बद्ध: जब तक भ्रंश सरकता रहता है,
नदी वहाँ तीव्र हो जाती है जहाँ वह उसे पार करती है। (ग) गतिशील: तट के
उत्थान ने हर नदी का \omterm{def:b3:quantitative-geomorphology:knickpoint}{आधार-तल} एक साथ नीचे कर दिया; \omterm{def:b3:quantitative-geomorphology:knickpoint}{निकपॉइंट}
ऊपरधारा की ओर खिसक रहे हैं, और गतिज-तरंग प्रमेय के अनुसार $n
= 1$ के लिए वे एक ही समोच्च-रेखा पर स्थित होते हैं।
\end{solution}

\begin{solution}{exo:b3:quantitative-geomorphology:5}
$\theta = m/n = 0.5$. $S = (U/K)^{1/2}A^{-1/2} = (10^{-4}/10^{-9})^{1/2}
\times(10^7)^{-1/2} = 316\times3.16\times10^{-4} = 0.10$। $U$ दोगुना करने पर
$k_s = (U/K)^{1/n}$ को $\sqrt2 = 1.41$ से गुणा करता है; $n = 1$ के लिए 2 से:
$n$ जितना बड़ा होगा, ढाल उतनी ही कम उत्थान पर प्रतिक्रिया देगी।
\end{solution}

\begin{solution}{exo:b3:quantitative-geomorphology:6}
निर्गम से ऊपरधारा की ओर $x$ लेने पर, $\mathrm{d}x = -\mathrm{d}d$, इसलिए $\chi(d)
= \int_d^L (A_0/c\,d'^{h})^{\theta}\,\mathrm{d}d' = A_0^{\theta} c^{-\theta}
\int_d^L d'^{-h\theta}\,\mathrm{d}d' = A_0^{\theta}c^{-\theta}(L^p -
d^p)/p$। यहाँ $p = 0.19$, $A_0^{\theta}c^{-\theta} = 10^{2.7}/2^{0.45}$,
और $\chi = \qty{2935}{m}$ प्राप्त होता है।
\end{solution}

\begin{solution}{exo:b3:quantitative-geomorphology:7}
$c = 2\times10^{-6}\times(10^8)^{0.5} = \qty{0.02}{m/yr}$, \qty{20}{m/kyr};
\qty{5}{km} इस गति से तय करने में \qty{250}{kyr} लगते हैं। क्षेत्रफल, और उसके साथ
चाल, ऊपरधारा की ओर घटती है, इसलिए \omterm{def:b3:quantitative-geomorphology:knickpoint}{निकपॉइंट} चढ़ते समय धीमा हो जाता है।
\end{solution}

\begin{solution}{exo:b3:quantitative-geomorphology:8}
$\bar e = 0.40\times600 = \qty{240}{m}$; शिखर $0.80\times240 =
\qty{192}{m}$ ऊपर उठते हैं; माध्य ऊँचाई $240 - 192 = \qty{48}{m}$ घटती है।
\end{solution}

\begin{solution}{exo:b3:quantitative-geomorphology:9}
$h = h_0\ln(P_0/E)$: $0.5\ln 6 = \qty{0.90}{m}$ और $0.5\ln1.5 =
\qty{0.20}{m}$। \qty{0.4}{mm/yr} पर अपरदन सबसे तेज़ उत्पादन,
$P_0$ से आगे निकल जाता है: मृदा उखड़ जाती है और ढाल उघड़ी चट्टान बन जाती है, जिसका निम्नीकरण
अपक्षय से सीमित रहता है।
\end{solution}

\begin{solution}{exo:b3:quantitative-geomorphology:10}
\omterm{def:b3:quantitative-geomorphology:steady-profile}{स्थायी अवस्था} में $\mathrm{d}q/\mathrm{d}x = E$ और जलविभाजक पर $q = 0$, इसलिए
$q = E x$। तब $\kappa S = E x\,(1 - S^2/S_c^2)$, एक द्विघात समीकरण जिसका धनात्मक
मूल $S = [-\kappa + \sqrt{\kappa^2 + 4E^2x^2/S_c^2}\,]/(2Ex/S_c^2)$ है।
20, 50, 100 और \qty{200}{m} पर: 0.36, 0.68, 0.89 और 1.03, जबकि 0.4, 1.0,
2.0 और 4.0 रैखिक विसरण के लिए ($S = Ex/\kappa$)। पहाड़ी का गोल
शिखर होता है, फिर ऐसी भुजाएँ जिनकी प्रवणता $S_c$ की ओर प्रवृत्त होती है: लगभग समतल ढाल,
परवलय की तुलना में कहीं कम उच्चावच।
\end{solution}

\begin{solution}{exo:b3:quantitative-geomorphology:11}
गतिज-तरंग प्रमेय के अनुसार परिवर्तन ऐसे बिंदु पर पहुँचता है जिसका निर्देशांक
$\chi$ है, $\chi/(KA_0^m)$ समय के बाद; नेटवर्क तब पूरी तरह प्रतिक्रिया दे चुका होता है जब
वह सबसे बड़े $\chi$ वाले अंतिम वाहिका-शीर्ष तक पहुँच चुका हो। $K A_0^m = 10^{-6}\times
1000 = \qty{e-3}{m/yr}$: मुख्य धारा \qty{5.46}{Myr} में प्रतिक्रिया देती है, नेटवर्क
\qty{7.24}{Myr} में, जो दस लाख वर्ष की अपेक्षित कोटि है।
\end{solution}

\begin{solution}{exo:b3:quantitative-geomorphology:12}
संगम के ऊपर सहायक नदी की स्थायी प्रवणता $S = (U/2K)A^{-\theta}$ है,
इसलिए $\mathrm{d}z/\mathrm{d}\chi = (U/2K)A_0^{-\theta}$, जो मुख्य धारा की आधी है।
\omterm{def:b3:quantitative-geomorphology:knickpoint}{निकपॉइंट} संगम तक तब पहुँचता है जब मुख्य धारा का $\chi_j$ पर होता है; उसके बाद एक और
समय $\Delta t$ के बाद मुख्य धारा $\chi_k$ तक खिसक चुकी है, जहाँ $\chi_k -
\chi_j = K A_0^m\,\Delta t$, जबकि सहायक नदी का \omterm{def:b3:quantitative-geomorphology:knickpoint}{निकपॉइंट}, चाल $2KA^m$ से,
$\chi$ तक खिसक चुका है, जहाँ $\chi - \chi_j = 2KA_0^m\,\Delta t = 2(\chi_k -
\chi_j)$।
\end{solution}

\begin{solution}{pb:b3:quantitative-geomorphology:1}
\textbf{1.} $4.5\times12000^{1.7} = \qty{3.87e7}{m^2}$, \qty{38.7}{km^2};
$4.5\times300^{1.7} = \qty{7.3e4}{m^2}$, \qty{0.073}{km^2}।

\textbf{2.} $\theta \approx -\log(0.007855/0.06641)/\log(37.08/0.388) =
0.928/1.980 = 0.47$।

\textbf{3.} हाँ: बारह बिंदुओं पर लगभग \qty{10}{\percent} प्रकीर्णन के साथ,
$0.466$ 0.5 से सार्थक रूप से भिन्न नहीं है।

\textbf{4.} $0.06641\times\sqrt{388\,000} = \qty{41.4}{m}$;
$0.007855\times\sqrt{3.708\times10^7} = \qty{47.8}{m}$। ज्यामितीय माध्य,
\qty{45.1}{m}, $k_s$ का $\theta = 0.5$ के लिए अनुमान देता है।

\textbf{5.} प्रवणता दो ऊँचाइयों का अंतर है, जिनमें से हर एक
डेसीमीटर से मीटर तक अनिश्चित है, एक छोटी दूरी से विभाजित; मंद भागों पर
त्रुटि अंतर का बड़ा भाग होती है। $\chi$ अवकलन के बजाय
समाकलन करता है।

\textbf{6.} $p = 1 - 1.7\times0.5 = 0.15$; $\chi = (1000/\sqrt{4.5})\,
(12000^{0.15} - 9000^{0.15})/0.15 = \qty{543.1}{m}$.

\textbf{7.} $(193.4 - 170)/543.1 = 0.0431$; $k_{sn} = 0.0431\times A_0^{0.5}
= \qty{43.1}{m}$.

\textbf{8.} काई अनुमान: यह नदी के \qty{3}{km} से अधिक भाग की
ऊँचाइयों पर टिका है, बारह शोरयुक्त प्रवणताओं पर नहीं।

\textbf{9.} $K = U/k_{sn} = 4.3\times10^{-5}/43.1 = \qty{1.00e-6}{yr^{-1}}$.

\textbf{10.} $10^{-6}\times\sqrt{3.87\times10^7}\times0.00691 =
\qty{4.30e-5}{m/yr}$, \qty{43}{m/Myr}: मुहाना उत्थान की
गति के साथ तालमेल बनाए रखता है।

\textbf{11.} मुहाने पर $10^{-6}\times6222 = \qty{6.2}{m/kyr}$;
जलप्रपात पर, $d = \qty{7.40}{km}$, $A = \qty{1.70e7}{m^2}$ और
\qty{4.1}{m/kyr}।

\textbf{12.} $170 + 43\times(4.5)^{-0.5}(12000^{0.15} - 300^{0.15})/0.15 =
\qty{405}{m}$: उच्चावच \qty{235}{m} है।

\textbf{13.} $n = 1$ के लिए समीकरण एक गतिक तरंग है: चरण के दोनों ओर के बिंदु
प्रति इकाई समय $U$ से ऊपर उठते हैं, इसलिए चरण अपनी ऊँचाई बनाए रखता है,
यानी \qty{25}{m} की गिरावट।

\textbf{14.} $\chi = (1000/\sqrt{4.5})(12000^{0.15} - 7400^{0.15})/0.15 =
\qty{899.4}{m}$.

\textbf{15.} $t = \chi/(KA_0^{0.5}) = 899.4/10^{-3} = \qty{899}{kyr}$.

\textbf{16.} $\chi = 899.4 + 300 = \qty{1199}{m}$ पर:
$d^{0.15} = 12000^{0.15} - 1199\times0.15\times\sqrt{4.5}/1000$, $d =
\qty{6.25}{km}$, मुहाने से \qty{5.75}{km} ऊपर की ओर।

\textbf{17.} संगम पर ($d = \qty{9.5}{km}$) $\chi = \qty{442.8}{m}$;
जलप्रपात वहाँ है जहाँ सहायक नदी का $\chi$ \qty{899.4}{m} तक पहुँचता है:
$s^{0.15} = 5000^{0.15} - (899.4 - 442.8)\times0.15\times\sqrt{1.5}/1000$,
$s = \qty{4.27}{km}$, संगम से \qty{0.73}{km} ऊपर।

\textbf{18.} उनके संगम $\chi = 1954$ और \qty{999}{m} पर हैं, जो
\qty{899}{m} से आगे हैं: \omterm{def:b3:quantitative-geomorphology:knickpoint}{निकपॉइंट} उन्हें अभी तक पार नहीं कर पाया है, इसलिए वहाँ मुख्य धारा
अब भी अपने पुराने स्तर पर है।

\textbf{19.} $E = 5\times10^{-3}\times0.0086 = \qty{4.3e-5}{m/yr}$,
\qty{43}{m/Myr}: जलप्रपात के ऊपर की पहाड़ियाँ उत्थान की दर से
अपरदित होती हैं, \omterm{def:b3:quantitative-geomorphology:steady-profile}{स्थायी अवस्था} में।

\textbf{20.} $150^2/(5\times10^{-3}) = \qty{4.5}{Myr}$, जो जलप्रपात की
आयु का पाँच गुना है।

\textbf{21.} उनका पाद वाहिका के साथ \qty{25}{m} नीचे हुआ: उनके
निचले भाग अधिक खड़े हो गए हैं, शायद सर्पण कर रहे हैं, जबकि उनके शिखर अब भी
पुरानी दर से अपरदित हो रहे हैं। प्रत्येक ढाल पर खड़ापन बढ़ने की एक तरंग चढ़ती है; कुछ
लाख वर्ष बाद वे उसी परवलय पर लौट आते हैं, \qty{25}{m} नीचे।

\textbf{22.} $t = \chi k_{sn}/(U A_0^{0.5})$ $1/U$ के अनुसार बदलता है:
$899/1.2 = 749$ से $899/0.8 = \qty{1124}{kyr}$ तक।

\textbf{23.} $n = 2$ पर वेग प्रवणता के साथ बढ़ता है: खड़ा
मुख आसपास के खंडों से तेज़ आगे बढ़ता और तीखा बना रहता, परंतु
उसकी गति और उससे पढ़ी गई आयु भिन्न होतीं।

\textbf{24.} वार्डन नदी का \omterm{def:b3:quantitative-geomorphology:knickpoint}{आधार-तल}
\textbf{\qty[text-series-to-math=true]{900}{kyr}} पहले गिरा (
परिकलन से \qty{899}{kyr})।
\end{solution}
